\documentclass[11pt]{article}
\usepackage[margin=1in]{geometry}
\usepackage{amsmath,amssymb,amsthm,mathtools}
\usepackage[T1]{fontenc}
\usepackage[utf8]{inputenc}
\usepackage{enumitem}
\usepackage{booktabs,longtable,array}
\usepackage{hyperref}
\usepackage{xcolor}
\usepackage{microtype}
\hypersetup{colorlinks=true, linkcolor=blue!50!black, citecolor=blue!50!black, urlcolor=blue!50!black}

\newtheorem{theorem}{Theorem}[section]
\newtheorem{proposition}[theorem]{Proposition}
\newtheorem{lemma}[theorem]{Lemma}
\newtheorem{corollary}[theorem]{Corollary}
\newtheorem{claim}[theorem]{Claim}
\theoremstyle{definition}
\newtheorem{definition}[theorem]{Definition}
\newtheorem{assumption}[theorem]{Assumption}
\theoremstyle{remark}
\newtheorem{remark}[theorem]{Remark}

\newcommand{\KL}{\mathrm{KL}}
\newcommand{\Ent}{\mathrm{Ent}}
\newcommand{\Prob}{\mathsf{P}}
\newcommand{\E}{\mathbb{E}}
\newcommand{\R}{\mathbb{R}}
\newcommand{\N}{\mathbb{N}}
\newcommand{\Pcal}{\mathcal{P}}
\newcommand{\Mcal}{\mathcal{M}}
\newcommand{\Qcal}{\mathcal{Q}}
\newcommand{\Zcal}{\mathcal{Z}}
\newcommand{\Fcal}{\mathcal{F}}
\newcommand{\Gcal}{\mathcal{G}}
\newcommand{\W}{\mathcal{W}}
\newcommand{\one}{\mathbf{1}}
\newcommand{\dd}{\,d}
\newcommand{\supp}{\operatorname{supp}}
\newcommand{\argmin}{\operatorname*{arg\,min}}
\newcommand{\argmax}{\operatorname*{arg\,max}}
\newcommand{\esssup}{\operatorname*{ess\,sup}}
\newcommand{\ri}{\operatorname{ri}}
\newcommand{\dom}{\operatorname{dom}}
\newcommand{\Law}{\operatorname{Law}}
\newcommand{\Var}{\operatorname{Var}}
\newcommand{\Cov}{\operatorname{Cov}}
\newcommand{\diag}{\operatorname{diag}}
\newcommand{\Id}{\mathrm{Id}}
\newcommand{\Lip}{\operatorname{Lip}}
\newcommand{\osc}{\operatorname{osc}}
\newcommand{\diam}{\operatorname{diam}}
\newcommand{\proj}{\operatorname{proj}}
\newcommand{\e}{\mathrm{e}}

\newenvironment{argument}{\paragraph{Argument.}\begin{enumerate}[leftmargin=2em]}{\end{enumerate}}
\newenvironment{proofoutline}{\paragraph{Proof architecture.}\begin{enumerate}[leftmargin=2em]}{\end{enumerate}}


% Source-faithfulness apparatus used by the paper reconstructions.
\newcommand{\sourceanchor}[1]{\par\smallskip\noindent\textbf{Source anchor.} #1\par\smallskip}
\newcommand{\sourcecomment}[1]{\begin{remark}#1\end{remark}}
\newenvironment{sourceguide}
  {\begin{description}[style=nextline,leftmargin=3.2cm,labelwidth=3cm]}
  {\end{description}}

% Backward-compatible environments used by some of the older reconstructions.
\newenvironment{distilled}[1][]{\begin{theorem}[#1]}{\end{theorem}}
\newenvironment{distillation}{\begin{enumerate}[leftmargin=2em]}{\end{enumerate}}

\newenvironment{commentarynotes}{\begin{remark}\begin{enumerate}[leftmargin=2em]}{\end{enumerate}\end{remark}}

\title{Variational Entropy, Convexity of Relative Entropy, and Supporting Slopes}
\author{}
\date{Classical results}
\begin{document}
\maketitle

\section*{Source map and scope}
\begin{description}[style=nextline,leftmargin=3.2cm,labelwidth=3cm]
\item[Primary source] Classical source results listed in the note.
\item[Coverage] A normalized proof dossier for the Donsker--Varadhan variational identity, Gibbs attainment, KL convexity/data processing, and one-dimensional supporting-slope arguments.
\item[Source anchors]
\begin{itemize}[leftmargin=2em]
  \item Donsker--Varadhan/Gibbs variational identity
  \item log-sum inequality
  \item data-processing variational proof
  \item supporting-line theorem for concave functions
\end{itemize}
\end{description}

The exposition below preserves the source proof order and proof technique.  Notation is normalized
and intermediate calculations are made explicit.  When the source contains a gap, ambiguity, or
apparently incorrect step, the reconstruction records the source step before adding a separate
commentary remark.

\section{Relative entropy and exponential tilting}
Let $(\Theta,\mathcal F)$ be a measurable space and let $\pi\in\mathcal P(\Theta)$.
For $\rho\in\mathcal P(\Theta)$ define
\[
 \KL(\rho\|\pi)=
 \begin{cases}
  \displaystyle\int\log\frac{d\rho}{d\pi}\,d\rho,&\rho\ll\pi,\\
  +\infty,&\text{otherwise}.
 \end{cases}
\]
Let $h:\Theta\to\mathbb R$ be measurable and suppose
$0<Z_h:=\int e^h\,d\pi<\infty$.  The exponentially tilted probability measure is
\begin{equation}\label{eq:CF-tilt}
 \pi_h(d\theta)=\frac{e^{h(\theta)}}{Z_h}\,\pi(d\theta).
\end{equation}
For every $\rho\ll\pi$ for which the terms are defined,
\begin{align}
 \KL(\rho\|\pi_h)
 &=\int\log\left(\frac{d\rho}{d\pi}e^{-h}Z_h\right)d\rho\\
 &=\KL(\rho\|\pi)-\int h\,d\rho+\log Z_h.
 \label{eq:CF-entropy-decomp}
\end{align}
Nonnegativity of relative entropy immediately yields the central change-of-measure
inequality.

\begin{theorem}[Donsker--Varadhan--Gibbs variational formula]
If $0<\int e^h\,d\pi<\infty$, then
\begin{equation}\label{eq:CF-DV}
 \log\int e^h\,d\pi
 =\sup_{\rho\in\mathcal P(\Theta)}
   \left\{\int h\,d\rho-\KL(\rho\|\pi)\right\}.
\end{equation}
The supremum is attained uniquely by $\pi_h$ whenever the objective is finite.
Equivalently, for every $\rho$,
\begin{equation}\label{eq:CF-change-measure}
 \int h\,d\rho\le\KL(\rho\|\pi)+\log\int e^h\,d\pi.
\end{equation}
\end{theorem}

\begin{proof}
Rearrange \eqref{eq:CF-entropy-decomp}:
\[
 \int h\,d\rho-\KL(\rho\|\pi)
 =\log Z_h-\KL(\rho\|\pi_h)\le\log Z_h.
\]
Equality holds for $\rho=\pi_h$, and only there by the equality condition in Gibbs'
inequality.
\end{proof}

The dual form is obtained by replacing $h$ with $-\lambda f$:
\begin{equation}\label{eq:CF-Gibbs-min}
 -\frac1\lambda\log\int e^{-\lambda f}\,d\pi
 =\inf_{\rho\in\mathcal P(\Theta)}
   \left\{\int f\,d\rho+\frac1\lambda\KL(\rho\|\pi)\right\},
\end{equation}
with minimizer proportional to $e^{-\lambda f}\pi$.

\section{Convexity of relative entropy}
Let $\rho_0,\rho_1,\pi_0,\pi_1$ be probability measures and $t\in[0,1]$.  Put
\[
 \rho_t=(1-t)\rho_0+t\rho_1,
 \qquad
 \pi_t=(1-t)\pi_0+t\pi_1.
\]
Choose a measure $m$ dominating all four measures and write their densities as
$r_i,p_i$.  The log-sum inequality states pointwise that
\begin{equation}\label{eq:CF-logsum}
 ((1-t)r_0+tr_1)
 \log\frac{(1-t)r_0+tr_1}{(1-t)p_0+tp_1}
 \le
 (1-t)r_0\log\frac{r_0}{p_0}+tr_1\log\frac{r_1}{p_1}.
\end{equation}
It follows from convexity of the perspective function
$(a,b)\mapsto a\log(a/b)$ on $\mathbb R_+^2$.  Integration yields:

\begin{theorem}[Joint convexity of relative entropy]
\begin{equation}\label{eq:CF-joint-convexity}
 \KL(\rho_t\|\pi_t)
 \le(1-t)\KL(\rho_0\|\pi_0)+t\KL(\rho_1\|\pi_1).
\end{equation}
In particular, $\rho\mapsto\KL(\rho\|\pi)$ is convex for fixed $\pi$.
\end{theorem}

\subsection{Mixtures with a common reference}
For fixed $\pi$, \eqref{eq:CF-joint-convexity} becomes
\begin{equation}\label{eq:CF-first-convexity}
 \KL((1-t)\rho_0+t\rho_1\|\pi)
 \le(1-t)\KL(\rho_0\|\pi)+t\KL(\rho_1\|\pi).
\end{equation}
This is the form used to prove concavity of ambiguity-set value functions: a mixture of
feasible distributions has no more than the corresponding mixture of their entropy
budgets.

\section{Data processing}
Let $T:\Theta\to\mathsf Y$ be measurable.  For bounded measurable $g$ on $\mathsf Y$,
apply \eqref{eq:CF-change-measure} to $h=g\circ T$:
\[
 \int g\,d(T_\#\rho)
 \le \KL(\rho\|\pi)+\log\int e^g\,d(T_\#\pi).
\]
Taking the supremum over $g$ and using \eqref{eq:CF-DV} on $\mathsf Y$ gives:

\begin{theorem}[Data-processing inequality]
\begin{equation}\label{eq:CF-DPI}
 \KL(T_\#\rho\|T_\#\pi)\le\KL(\rho\|\pi).
\end{equation}
\end{theorem}

For a Markov kernel $K$ from $\Theta$ to $\mathsf Y$, the same argument on the joint
measures $\rho(d\theta)K(\theta,dy)$ and $\pi(d\theta)K(\theta,dy)$ yields
\[
 \KL(\rho K\|\pi K)\le\KL(\rho\|\pi).
\]

\section{Conditional entropy chain rule}
Let $P$ and $Q$ be probability measures on $X\times Y$ with disintegrations
\[
 P(dx,dy)=P_X(dx)P^x(dy),
 \qquad
 Q(dx,dy)=Q_X(dx)Q^x(dy).
\]
Whenever $P\ll Q$, Radon--Nikodym factorization gives
\[
 \frac{dP}{dQ}(x,y)
 =\frac{dP_X}{dQ_X}(x)\frac{dP^x}{dQ^x}(y)
\quad P\text{-almost everywhere}.
\]
Taking logarithms and integrating gives:
\begin{theorem}[Entropy chain rule]
\begin{equation}\label{eq:CF-chain-rule}
 \KL(P\|Q)=\KL(P_X\|Q_X)
 +\int_X\KL(P^x\|Q^x)\,P_X(dx).
\end{equation}
\end{theorem}
Both terms are nonnegative.  Equality in data processing for the projection
$(x,y)\mapsto x$ occurs precisely when $P^x=Q^x$ for $P_X$-almost every $x$.

\section{Concave value functions}
Let $F$ be a linear functional on probability measures and let $C$ be a nonnegative
convex constraint functional.  Define
\begin{equation}\label{eq:CF-value}
 v(r)=\sup\{F(\rho):C(\rho)\le r\},
 \qquad r\ge0.
\end{equation}
If $\rho_i$ is feasible at radius $r_i$, then
$\rho_t=(1-t)\rho_0+t\rho_1$ satisfies
\[
 C(\rho_t)\le(1-t)C(\rho_0)+tC(\rho_1)
 \le(1-t)r_0+tr_1,
\]
and linearity gives
$F(\rho_t)=(1-t)F(\rho_0)+tF(\rho_1)$.  Taking suprema proves:

\begin{proposition}[Concavity and monotonicity]
The value function $v$ in \eqref{eq:CF-value} is nondecreasing and concave on its
effective domain.
\end{proposition}

\section{Supporting slopes for one-dimensional concave functions}
Let $v:[a,b]\to\mathbb R$ be concave and let $r\in(a,b)$.  For $x<r<y$, concavity
implies the monotonicity of secant slopes:
\begin{equation}\label{eq:CF-secant-order}
 \frac{v(r)-v(x)}{r-x}
 \ge
 \frac{v(y)-v(x)}{y-x}
 \ge
 \frac{v(y)-v(r)}{y-r}.
\end{equation}
Define the one-sided slope bounds
\begin{align}
 D^-v(r)&=\inf_{x<r}\frac{v(r)-v(x)}{r-x},\\
 D^+v(r)&=\sup_{y>r}\frac{v(y)-v(r)}{y-r}.
\end{align}
Then $D^+v(r)\le D^-v(r)$.  Every
$\lambda\in[D^+v(r),D^-v(r)]$ satisfies
\begin{equation}\label{eq:CF-supergradient}
 v(s)\le v(r)+\lambda(s-r)
 \qquad(s\in[a,b]).
\end{equation}
Such a $\lambda$ is a supergradient of $v$ at $r$.

\begin{proposition}[Nonnegative supergradient]
If $v$ is also nondecreasing, every supergradient chosen from the nonnegative portion of
$[D^+v(r),D^-v(r)]$ is nonnegative.  In particular, at every interior point of a finite,
nondecreasing concave value function there exists $\lambda\ge0$ satisfying
\eqref{eq:CF-supergradient}.
\end{proposition}

\begin{proof}
For $y>r$, monotonicity gives
$(v(y)-v(r))/(y-r)\ge0$, hence $D^+v(r)\ge0$.  The interval of admissible supporting
slopes is nonempty by \eqref{eq:CF-secant-order}, so it contains a nonnegative slope.
\end{proof}

\section{Lagrange multiplier consequence}
Apply \eqref{eq:CF-supergradient} to the value function \eqref{eq:CF-value}.  If
$C(\rho)=s$, then $F(\rho)\le v(s)$, and therefore
\[
 F(\rho)\le v(r)+\lambda(C(\rho)-r).
\]
Rearranging gives the global Lagrangian bound
\begin{equation}\label{eq:CF-lagrangian}
 F(\rho)-\lambda C(\rho)\le v(r)-\lambda r.
\end{equation}
Taking a supremum over $\rho$ yields
\begin{equation}\label{eq:CF-dual-value}
 \sup_\rho\{F(\rho)-\lambda C(\rho)\}+\lambda r\le v(r).
\end{equation}
The reverse inequality follows from weak duality for every $\lambda\ge0$.  Thus the
supporting slope produces an optimal Lagrange multiplier whenever the inner supremum
can be evaluated or approximated.

\section{How the pieces combine}
The recurring argument is:
\begin{enumerate}
\item convexity of entropy or transport cost makes the budget value function concave;
\item one-dimensional concavity produces a nonnegative supporting slope;
\item measurable near-maximizers construct a distribution whose Lagrangian value
approaches the inner supremum;
\item the supporting-line inequality and weak duality meet at equality;
\item the Donsker--Varadhan formula evaluates entropy-penalized inner problems by an
exponential tilt.
\end{enumerate}
These are independent classical ingredients, but together they form a general mechanism
for strong duality and Gibbs-type optimizer formulas.

\begin{thebibliography}{99}
\bibitem{ChenGeorgiouPavon2021} Yongxin Chen, Tryphon T. Georgiou, and Michele Pavon. \emph{Stochastic Control Liaisons: Richard Sinkhorn Meets Gaspard Monge on a Schr\"odinger Bridge}. SIAM Review, 63(2):249--313, 2021. doi:10.1137/20M1339982. arXiv:2005.10963.
\bibitem{Leonard2014Survey} Christian L\'eonard. \emph{A Survey of the Schr\"odinger Problem and Some of Its Connections with Optimal Transport}. Discrete and Continuous Dynamical Systems--A, 34(4):1533--1574, 2014. doi:10.3934/dcds.2014.34.1533. arXiv:1308.0215.
\bibitem{Fortet1940} Robert Fortet. \emph{Resolution d'un systeme d'equations de M. Schrodinger}. Journal de mathematiques pures et appliquees, 9e serie, 19(1--4):83--105, 1940.
\bibitem{SinkhornKnopp1967} Richard Sinkhorn and Paul Knopp. \emph{Concerning Nonnegative Matrices and Doubly Stochastic Matrices}. Pacific Journal of Mathematics, 21(2):343--348, 1967.
\bibitem{FranklinLorenz1989} Joel Franklin and Jens Lorenz. \emph{On the Scaling of Multidimensional Matrices}. Linear Algebra and its Applications, 114/115:717--735, 1989. doi:10.1016/0024-3795(89)90490-4.
\bibitem{Carroll2004} Joseph E. Carroll. \emph{Birkhoff's Contraction Coefficient}. Linear Algebra and its Applications, 389:227--234, 2004. doi:10.1016/j.laa.2004.02.039.
\bibitem{EvesonNussbaum1995} Simon P. Eveson and Roger D. Nussbaum. \emph{An Elementary Proof of the Birkhoff--Hopf Theorem}. Mathematical Proceedings of the Cambridge Philosophical Society, 117(1):31--55, 1995.
\bibitem{BlanchetMurthy2019} Jose Blanchet and Karthyek R. A. Murthy. \emph{Quantifying Distributional Model Risk via Optimal Transport}. Mathematics of Operations Research, 44(2):565--600, 2019. doi:10.1287/moor.2018.0936. arXiv:1604.01446.
\bibitem{GaoKleywegt2023} Rui Gao and Anton J. Kleywegt. \emph{Distributionally Robust Stochastic Optimization with Wasserstein Distance}. Mathematics of Operations Research, 48(2):603--655, 2023. doi:10.1287/moor.2022.1275. arXiv:1604.02199.
\bibitem{MohajerinEsfahaniKuhn2018} Peyman Mohajerin Esfahani and Daniel Kuhn. \emph{Data-Driven Distributionally Robust Optimization Using the Wasserstein Metric: Performance Guarantees and Tractable Reformulations}. Mathematical Programming, 171(1--2):115--166, 2018. doi:10.1007/s10107-017-1172-1. arXiv:1505.05116.
\bibitem{WangGaoXie2025} Jie Wang, Rui Gao, and Yao Xie. \emph{Sinkhorn Distributionally Robust Optimization}. Operations Research, 2025. doi:10.1287/opre.2023.0294. arXiv:2109.11926.
\bibitem{Girsanov1960} I. V. Girsanov. \emph{On Transforming a Certain Class of Stochastic Processes by Absolutely Continuous Substitution of Measures}. Theory of Probability and Its Applications, 5(3):285--301, 1960. doi:10.1137/1105027.
\bibitem{CameronMartin1945} R. H. Cameron and W. T. Martin. \emph{Transformations of Wiener Integrals under a General Class of Linear Transformations}. Transactions of the American Mathematical Society, 58(2):184--219, 1945.
\bibitem{Yeh1978} J. Yeh. \emph{Transformation of Conditional Wiener Integrals under Translation and the Cameron--Martin Translation Theorem}. Tohoku Mathematical Journal, 30(4):505--515, 1978.
\bibitem{Kakutani1948} Shizuo Kakutani. \emph{On Equivalence of Infinite Product Measures}. Annals of Mathematics, 49(1):214--224, 1948.
\bibitem{ShiDeBortoliCampbellDoucet2023} Yuyang Shi, Valentin De Bortoli, Andrew Campbell, and Arnaud Doucet. \emph{Diffusion Schr\"odinger Bridge Matching}. Advances in Neural Information Processing Systems, 2023. arXiv:2303.16852.
\bibitem{LiuLipmanNickelKarrerTheodorouChen2024} Guan-Horng Liu, Yaron Lipman, Maximilian Nickel, Brian Karrer, Evangelos A. Theodorou, and Ricky T. Q. Chen. \emph{Generalized Schr\"odinger Bridge Matching}. International Conference on Learning Representations, 2024. arXiv:2310.02233.
\bibitem{DomingoEnrichDrozdzalKarrerChen2025} Carles Domingo-Enrich, Michal Drozdzal, Brian Karrer, and Ricky T. Q. Chen. \emph{Adjoint Matching: Fine-Tuning Flow and Diffusion Generative Models with Memoryless Stochastic Optimal Control}. arXiv preprint arXiv:2409.08861, 2025.

\bibitem{Birkhoff1957} Garrett Birkhoff. \emph{Extensions of Jentzsch's Theorem}. Transactions of the American Mathematical Society, 85(1):219--227, 1957.
\bibitem{Sinkhorn1964} Richard Sinkhorn. \emph{A Relationship Between Arbitrary Positive Matrices and Doubly Stochastic Matrices}. Annals of Mathematical Statistics, 35(2):876--879, 1964.
\bibitem{Sinkhorn1967Prescribed} Richard Sinkhorn. \emph{Diagonal Equivalence to Matrices with Prescribed Row and Column Sums}. American Mathematical Monthly, 74(4):402--405, 1967.
\bibitem{Bauer1965} Friedrich L. Bauer. \emph{An Elementary Proof of the Hopf Inequality for Positive Operators}. Numerische Mathematik, 7:331--337, 1965.
\bibitem{Bushell1973} P. J. Bushell. \emph{Hilbert's Metric and Positive Contraction Mappings in a Banach Space}. Archive for Rational Mechanics and Analysis, 52:330--338, 1973.

\end{thebibliography}

\end{document}
